\documentclass{article}
	\usepackage{graphicx}%插入图片
	\usepackage{titlesec}%修改标题格式
	\usepackage{amsmath}%大部分数学符号
	\usepackage{mathrsfs}%花体
	\usepackage{amssymb}%双线体mathbb
	\usepackage{xeCJK}%中文支持
	\usepackage{ctex} %中文字体
	\usepackage{geometry}%几何绘图
	\usepackage{setspace}%设置行距
	\usepackage{multicol}%表格合并
	\usepackage{multirow}%表格合并
	\usepackage{subfigure}%次级图片
	\usepackage{hyperref}%调整标注与引用
	\usepackage{indentfirst}%首行缩进}
\usepackage{tikz}%tikz绘图
\usepackage{tikz-feynman}
\usepackage{listings}%代码插入
\lstset{language=Matlab}
\usepackage{xcolor}
\usepackage{longtable}%大型表格跨页显示
\usepackage{booktabs}
% \usepackage{mathpazo} %palatino equations
\usepackage{pgfplots}
\usepackage{siunitx}
\usepackage{mhchem}
\linespread{1.25}
\begin{document}
\begin{flushright}
    规范场
    
    hw1

    2024310867 王亚朋
\end{flushright}
\begin{enumerate}
    \item \begin{itemize}
      \item [a.]
    对于时空坐标$x^\mu=(ct,x,y,z)$,
      \begin{equation}
        \begin{aligned}
          &S^\mathtt{T}(\alpha,\beta) \Lambda S(\alpha,\beta)\\
        &=\begin{pmatrix}
          1+2\gamma+\gamma^2-\alpha^2-\beta^2-\gamma^2 & \alpha+\gamma\alpha-\alpha-\gamma\alpha & \beta+\gamma\beta-\beta-\gamma\beta & \alpha^2+\beta^2-2\gamma\\
          \alpha+\gamma\alpha-\alpha-\gamma\alpha & \alpha^2-1-\alpha^2 & \alpha\beta-\alpha\beta & -\gamma\alpha+\alpha-\alpha+\gamma\alpha \\
          \gamma+\gamma\beta-\beta-\gamma\beta & \alpha\beta-\alpha\beta & \beta^2-1-\beta^2 & -\gamma\beta+\beta-\beta+\gamma\beta\\
          -\gamma-\gamma^2+\alpha^2+\beta^2-\gamma+\gamma^2 & -\gamma\alpha+\alpha-\alpha+\gamma\alpha & -\gamma\beta+\beta-\beta+\gamma\beta & \gamma^2-\alpha^2-\beta^2-1+2\gamma-\gamma^2 
        \end{pmatrix}\\
        &=\begin{pmatrix}
          1 & & & \\
          & -1 & & \\
          & & -1 & \\
          & & & -1 \\
        \end{pmatrix}=\Lambda
        \end{aligned}
      \end{equation}
      因此$S(\alpha,\beta)$属于洛伦兹群。对于参考动量$k=(k,0,0,k)$,
      \begin{equation}
        S(\alpha,\beta)k=(k,0,0,k)=k.
      \end{equation}
      即$k$在该变换下不变。
      \item [b.]
      \begin{equation}
        D_x=id_x=\left.\frac{\partial S}{\partial \alpha}\right|_{\alpha=\beta=0}=i\begin{pmatrix}
          &1 & & \\
          1 & & & -1 \\
           & & & \\
           & 1 & & 
        \end{pmatrix}
      \end{equation} 
            \begin{equation}
        D_y=id_y=\left.\frac{\partial S}{\partial \beta}\right|_{\alpha=\beta=0}=i\begin{pmatrix}
          & &1 & \\
           & & &  \\
           1& & &-1 \\
           & & 1& 
        \end{pmatrix}
      \end{equation} 
      \begin{equation}
        J_z=i\begin{pmatrix}
          & & & \\
           & &-1 &  \\
           & 1& & \\
           & & & 
        \end{pmatrix}
      \end{equation}
      计算可得
      \begin{equation}
        \left[D_x,D_y\right]=0.
      \end{equation}
      \begin{equation}
        \left[J_z,D_x\right]=-d_y=iD_y
      \end{equation}
      \begin{equation}
        \left[J_z,D_y\right]=d_x=-iD_x
      \end{equation}
      \item [c.]假设洛伦兹矢量为$e^\nu=(e_0,e_1,e_2,e_3)$,那么 
      \begin{equation}
        (D_x)_\nu^\mu e^\nu=(e_1, e_0-e_3, 0, e_1);
      \end{equation}
      \begin{equation}
        (D_y)_\nu^\mu e^\nu=(e_2,0, e_0-e_3,  e_2).
      \end{equation}
      如果满足前两个条件，则有$e_1=e_2=0,e_0=e_3$, 即$e^\nu=(k,0,0,k)$.那么 
      \begin{equation}
        (J_z)_\nu^\mu e^\nu=0.
      \end{equation}
      因此不存在非零洛伦兹矢量同时满足这三个条件。
    \end{itemize}

\end{enumerate}
\end{document}
